%% ============================================================
%% 分片 12：原文第 5 节「Topological Field Theory In Two Dimensions」（原文 glshort.tex 起于第 3264 行）
%% 由 tools/harvmac2tex.py 从 arXiv 源码（hep-th/0604151v3）自动生成。
%% 只译英文散文；公式、编号、\label、\cite 一律不改。
%% ============================================================

\subsection{扭化拓扑场论}

\label{sec:twistedtft}

我们现在想从拓扑场论的观点来研究靶空间为 $\MH$ 的 sigma 模型，因此把
$\Bbb{R}^{1,1}$ 换成一个一般的二维流形 $\Sigma$，最终我们将取它带有欧几里得号差。我们想
讨论：由靶空间为 $\MH$ 的 sigma 模型可以构造出 $\Sigma$ 上的哪些扭化拓扑场论，以及其中哪些实际上来自把我们在第~\ref{sec:tqftfour}~节中
构造的四维拓扑场论在黎曼面 $C$ 上紧化。

 \cite{ref66}
 \cite{ref67}
\cite{ref68}


为了由靶空间为 $X$ 的 sigma 模型构造拓扑场论，人们在 $X$ 上取一\emph{对}满足 $(2,2)$ 超对称某些条件的
复结构 $(J_+,J_-)$ 。\footnote{广义复几何给出这一构造的一个推广，我们这里不需要它：在推广中，由对 $(J_+,J_-)$ 可以形成的两个广义复结构中 \cite{ref67}\ 只有一个可积。
然后人们像  那样，试图用这个可积的广义复结构来定义拓扑场论。}这些
条件最近被重新解释 ，所用语言是广义复几何 \cite{ref69}。一旦有了 $(2,2)$ 超对称，就有一个
标准配方 \cite{ref70}（通过我们在第~\ref{sec:twisting}~节中回顾过的扭化程序）来构造拓扑
场论。

若 $X$ 是超凯勒的，事情会简化。靶空间为 $X$ 的
sigma 模型具有 $(4,4)$ 超对称，而 $(2,2)$ 超对称的结构可以通过取
$(4,4)$ 超对称代数的一个 $(2,2)$ 子代数得到。\footnote{相差 $Spin(4)_{\cal R}$ 共轭的两个 $(2,2)$ 子代数
视为等价的。}就复结构的选择而言，这相当于以下内容。$X$ 有一族由一份
$\Bbb{CP}^1$ 参数化的复结构，我们把它记作
$\Bbb{CP}^1_h$；而取一对点 $(J_+,J_-)\in \Bbb{CP}_h^1$ 就可以定义一个 $(2,2)$ 结构。\cite{ref66}\ 中的条件对这样一对点
自动满足。对 $\MH$ 而言，这就给出对几何 Langlands 纲领重要的那族
拓扑场论。

我们可以如下方便地刻画与对 $(J_+,J_-)$ 相联的拓扑场论。设 $\Phi: \Sigma\to
\MH$ 是 sigma 模型的玻色场。那么未破缺超对称的方程为
% 式 (5.5)
\setcounter{equation}{4}
\begin{equation}\label{eq:hson}
(1-iJ_+)\bar\partial\Phi=0,
\end{equation}
 以及
% 式 (5.6)
\setcounter{equation}{5}
\begin{equation}\label{eq:hsono}
(1-iJ_-)\partial \Phi=0.
\end{equation}
  这里 $\partial$ 与
$\bar\partial$ 是 $\Sigma$ 上外微分的通常 $(1,0)$ 与 $(0,1)$ 部分；而 $\partial\Phi$ 与
$\bar\partial \Phi$ 理解为 $\Sigma$ 上取值于切丛 $TX$ 的拉回（$J_+$
与 $J_-$ 作用在其上）的 1-形式。第一个条件说映射
$\Phi:\Sigma\to X$ 在复结构 $J_+$ 中全纯；第二个说它在复结构 $J_-$ 中反全纯。
与我们在第~\ref{sec:familya}~节中找出未破缺超对称条件的方式类似，这些方程之所以出现，是因为存在费米场 $\chi_+$、$\chi_-$，使得
$\{Q,\chi_+\}=(1-iJ_+)\bar\partial\Phi$、
$\{Q,\chi_-\}=(1-iJ_-)\partial\Phi$。

方程 \eqref{eq:hson}\ 与 \eqref{eq:hsono}\ 中的信息远多于其解中的信息。为了讨论
方程的解，人们通常特殊化到欧几里得号差的 $\Sigma$（对拓扑场论来说这是最自然的情形）。
此时对一般的 $J_+$、$J_-$，它们蕴含映射
$\Phi:\Sigma\to X$ 必须是常值映射，而这个结果并不依赖于
$J_+$ 与 $J_-$。方程本身若从代数上理解、并且不固定二维流形
$\Sigma$ 上的实结构，则依赖于 $J_+$ 与 $J_-$。

一个重要的例子是 $J_+=J_-=\hat J$ 的情形（在我们的应用中，$\hat J$ 将是 $\MH$ 的复结构 $I_w$
之一，如式 \eqref{eq:wyre} 所述）。这些方程合起来给出
$(1-i\hat J\,)d\Phi = 0$，从而蕴含 $d\Phi=0$。该模型
称为复结构 $\hat J$ 中的 $B$-模型。反之若
$J_+=-J_-=\hat J$，我们就得到复结构 $\hat J$ 中的 $A$-模型。（这个模型实际上依赖于与复结构 $\hat J$ 相联的辛形式
$\omega_{\hat J}$，而不是复结构本身。但有时把「复结构 $\hat J$ 中的 $A$-模型」当作方便的简称。）对欧几里得号差的 $\Sigma$，
$A$-模型的方程是多余的，因为式 \eqref{eq:hsono}\ 是 \eqref{eq:hson}\ 的复共轭。
其中任何一个都断言映射 $\Phi:\Sigma\to \MH$ 在复结构 $\hat
J$ 中全纯。若 $J_-\not= \pm J_+$，我们得到的模型既不是
$A$-模型也不是 $B$-模型，但可以用广义复几何把它约化为 $A$-模型，正如我们在第~\ref{sec:gencpx}~节中讨论的那样。

我们回顾式 \eqref{eq:wyre}\ 中由 $\Bbb{CP}_h^1$ 参数化的 $\MH$ 上复结构族的显式形式。用
$\Bbb{CP}_h^1$ 的仿射参数 $w$ 来说，我们定义了一个复结构 $I_w$，其全纯坐标是
$A_{\bar z}-w\phi_{\bar z}$ 与 $A_z+w^{-1}\phi_z$。具体地，我们
得到 % 式 (5.7)
\setcounter{equation}{6}
\begin{equation}\label{eq:wyrre}
I_w={1-\bar w w\over 1+\bar w w}I+{i(w-\bar
w)\over 1+\bar w w}J+{w+\bar w\over 1+\bar w w}K.
\end{equation}

 在这个参数化下，对
$(J_+,J_-)$ 对应于对
$(w_+,w_-)$。

\par\medskip\emph{从四维约化}

现在让我们把它与从四维约化所得的结果作比较。我们在第~\ref{sec:tqftfour}~节中了解到，规范理论
给出定义在任何四维流形 $M$ 上、并由一份我们记作
$\Bbb{CP}^1_g$ 的 $\Bbb{CP}^1$ 参数化的一族拓扑量子场论（TQFT）。
这一族对 $M=\Sigma\times C$ 当然约化为族
$\Bbb{CP}_h^1\times \Bbb{CP}_h^1$ 中从二维观点看自然的子族。两族都来自对
四维流形 $M$ 上 ${\EUN}=4$ 超对称杨-米尔斯理论施行扭化程序。在该程序中，拓扑场论
由适当选取的超对称线性组合 $Q$ 的上同调得到。要得到族 $\Bbb{CP}^1_g$，我们
要求 $Q$ 是 $Spin'(4)$-不变的，以便构造对所有的 $M$ 都有效。但要得到靶空间为 $\MH$ 的 sigma 模型，我们
特殊化到约化和乐流形 $M=\Sigma\times C$，此时如我们后面更明确讨论的，对 $Q$ 的较弱条件就够了，从而得到更大的族 $\Bbb{CP}_h^1\times
\Bbb{CP}_h^1$。

我们希望把四维拓扑场论族 $\Bbb{CP}^1_g$ 辨认作 sigma 模型族
$\Bbb{CP}_h^1\times \Bbb{CP}_h^1$ 的子族。$\Bbb{CP}^1_g$
到 $\Bbb{CP}_h^1\times \Bbb{CP}_h^1$ 的嵌入可以用函数
$w_+(t)$、$w_-(t)$ 描述。通过特殊化到一个方便的构型或物理态，我们可以算出这些函数。

一个非常方便的做法是干脆把问题阿贝尔化，即在一个由某些未扭化标量场的
期望值把 $G$ 破缺到其极大环面的真空中工作。在阿贝尔情形，由 \eqref{eq:loopme}，四维中未破缺超对称的方程为
% 式 (5.8)
\setcounter{equation}{7}
\begin{equation}\label{eq:onso}
\begin{aligned}(F+t\,d\phi)^+&=0\\
(F-t^{-1}d\phi)^-&=0.\\
\end{aligned}
\end{equation}

这里，1-形式 $\phi$ 是 $T^*M=T^*\Sigma\oplus
T^*C$ 的截面，而规范场 $A$ 局部地是一个 1-形式或 $T^*M$ 的截面。为了与 $\Sigma$ 上靶空间为
$\MH(G,C)$ 的 sigma 模型比较，我们必须取 $\Sigma $ 的半径远大于
$C$ 的半径。在该极限下，$A$ 与 $\phi$ 中取值于
$T^*\Sigma$ 的部分是「有质量的」，可以丢掉；于是
$A$ 与 $\phi$ 可以解释为 $T^*C$ 的截面，在 $\Sigma$ 上缓慢
变化。

我们用 $y$ 记 $\Sigma$ 上的局部复坐标，用 $z$
记 $C$ 上的局部复坐标。2-形式 $d\bar
y\wedge d\bar z$ 是自对偶的，而 $dy\wedge d\bar z$ 是
反自对偶的。因此 \eqref{eq:onso} 中的第一个方程给出
% 式 (5.9)
\setcounter{equation}{8}
\begin{equation}\label{eq:igot}
\partial_{\bar y}\left(A_{\bar z}+t \phi_{\bar
z}\right)=0,
\end{equation}
 而第二个方程给出
% 式 (5.10)
\setcounter{equation}{9}
\begin{equation}\label{eq:nigot}
\partial_y\left(A_{\bar z}-t^{-1}\phi_{\bar
z}\right)=0.
\end{equation}

把这些结果与式 \eqref{eq:hson}\ 和 \eqref{eq:hsono}\ 比较，就可以读出函数 $w_+(t)$ 与 $w_-(t)$。式 \eqref{eq:igot}\ 断言：
若在 $\MH(G,C)$ 上取使 $A_{\bar z}+t\phi_{\bar
z}$ 全纯的复结构，则映射 $\Sigma\to \MH(G,C)$ 全纯。这与式 \eqref{eq:hson} 一致；后者断言
在复结构 $I_{w_+}$ 中全纯，当且仅当
$w_+=-t$。类似地，式 \eqref{eq:nigot}\ 断言映射 $\Sigma\to
\MH(G,C)$ 反全纯，如果我们在 $\MH(G,C)$ 上取使 $A_{\bar z}-t^{-1} \phi_{\bar z}$
全纯的复结构。这与式 \eqref{eq:hsono} 一致；后者断言在复结构 $I_{w_-}$ 中
反全纯，当且仅当
$w_-=t^{-1}$。

于是 $\Bbb{CP}^1_g$ 到 $\Bbb{CP}_h^1\times
\Bbb{CP}_h^1$ 的嵌入由 % 式 (5.11)
\setcounter{equation}{10}
\begin{equation}\label{eq:tyroy}
\begin{aligned}w_+ & = -t\\
                                              w_- & = t^{-1}.\\
\end{aligned}
\end{equation}
定义。

\par\medskip\emph{一些简单的考虑}

现在让我们对这个结果作几点简单的观察。

首先，我们什么时候得到 $B$-模型？对 $B$-模型，我们
要求 $w_+=w_-$。所以条件是 $t^{-1}=-t$，它恰好出现在
$t=\pm i$ 时。由于 $w=\mp i$ 时的复结构 $I_w$ 与 $J$、$-J$ 重合，这样我们就得到
复结构 $J$ 或 $-J$ 中的 $B$-模型。

我们什么时候得到 $A$-模型？对 $A$-模型，复结构
$I_{w_+}$ 与 $I_{w_-}$ 应当相反。映射
$w\to -1/\bar w$ 把 $I_w$ 映到它的相反结构（这从
\eqref{eq:wyrre} 看得很清楚），所以若 $w_-=-1/\bar
w_+$，我们就得到二维的 $A$-模型，这归结为 $t=\bar t$。换句话说，恰对实的 $t$，我们得到 $A$-模型。
例如对 $t=1$ 或
$-1$，我们得到复结构 $K$ 或 $-K$ 中的 $A$-模型。对
$t=0$ 或 $\infty$，我们得到复结构 $I$
或 $-I$ 中的 $ A$-模型。

现在，让我们把这个族与 sigma 模型更完整的族
$\Bbb{CP}_h^1\times \Bbb{CP}_h^1$ 作比较。在第~\ref{sec:complexstr}~节末尾，我们描述了 $\MH(G,C)$ 的微分同胚群 ${\cal U}\cong
\Bbb{C}^*$；它作用在复结构族 $I_w$ 上为 $w\to \lambda^{-1} w$，$\lambda\in
\Bbb{C}^*$。由对\footnote{这个拓扑场论还依赖于 $\MH$ 的超凯勒度量与 $B$-场。
对全部变量的完整依赖关系将在第~\ref{sec:gencpx}~节中用广义复几何确定。}
$(w_+,w_-)$ 确定的拓扑场论与由 $(\lambda^{-1} w_+,\lambda^{-1} w_-)$ 确定的相同。从二维观点看，我们能定义多少种真正不等价的 TQFT？
由对 $(w_+,w_-)$ 能构成的唯一不变量是比值 $q=w_+/w_-$。我们看到
\begin{equation*}
\tag*{(5.12)}
q=-t^2,
\end{equation*}
 所以 $q$ 的所有值都可以实现，但并非完全唯一。$\Bbb{CP}^1_g$ 上的点 $t$ 与 $-t $ 导致二维中相互等价的理论。这一等价
反映了 $SU(4)_{\cal R}$ 的中心的作用（回忆式
\eqref{eq:ungu}）。

尽管不变量 $q$ 的所有值确实都来自
四维 TQFT，但并非 $\Bbb{CP}_h^1\times \Bbb{CP}_h^1$ 上所有的 $\Bbb{C}^*$
轨道都有这种来源的代表元。缺失的轨道是
$\Bbb{C}^*$-不变点 $(0,0)$ 与 $(\infty,\infty)$，以及
$w_+$ 或 $w_-$（而不是两者）为 0 或
$\infty$ 的轨道。特别值得注意的是，点 $(0,0)$
与 $(\infty,\infty)$ 并不等价于源于四维 TQFT 的理论。这些点对应于
复结构 $I$ 与 $-I$ 中的 $B$-模型。

复结构 $I$ 中的 $B$-模型一直是数学方面努力的出发点
——引言中已简要概述过——这些努力要把几何 Langlands 纲领解释为
$\MH$ 的几何。由于 Hitchin 纤维化在复结构 $I$ 中全纯，Hitchin 纤维化纤维上的
$T$-对偶（它与四维 $S$-对偶的关系我们很快就会回顾）把复结构 $I$ 的 $B$-模型映到
自身，并通过 Fourier-Mukai 变换作用在 $D$-膜上，而后者正是数学描述的出发点。尽管
与这个 $B$-模型对应的点 $(0,0)$ 不在来自四维 TQFT 的族 $\Bbb{CP}_g^1$ 中，但
有趣的是，它可以被无穷小地扰动，从而给出与 $\Bbb{CP}^1_g$ 上几乎任意点对应的 $\Bbb{C}^*$ 轨道上的点。
我们只需把 $(0,0)$ 扰动到 $(\alpha,\beta)$，其中 $\alpha$、$\beta$ 无穷小；不变量 $q=w_+/w_-$ 于是为 $q=\alpha/\beta$，而对任意小的
$\alpha,\beta$ 它可以取任何值。也许这一事实最终会带来一种基于从复结构 $I$ 的 $B$-模型出发作扰动的几何 Langlands 纲领理解。
然而，我们的方法将依赖于直接来自四维的族 $\Bbb{CP}^1_g$。

尽管某些 $\Bbb{C}^*$ 轨道（例如点 $(0,0)$）不能通过把四维 TQFT 特殊化到
$M=\Sigma\times C$ 而出现，这并不意味着它们不能用四维规范理论描述。在第~\ref{sec:twisting}~节中，
我们利用被约束为 $Spin'(4)$-不变的超对称生成元 % 式 (5.13)
\setcounter{equation}{12}
\begin{equation}\label{eq:klop}
\epsilon=u\epsilon_\ell+v\epsilon_\rr
\end{equation}
 得到了族 $\Bbb{CP}^1_g$。
若特殊化到
$M=\Sigma\times C$，就不必要求 $Spin'(4)$
不变性。$\Sigma\times C$ 的和乐群是
$SO(2)\times SO(2)\subset SO(4)$，而 $Spin'(2)\times Spin'(2)$
不变性就足够了。这意味着我们可以把 \eqref{eq:klop}\ 推广为
% 式 (5.14)
\setcounter{equation}{13}
\begin{equation}\label{eq:blop}
\hat\epsilon=\left(u+\tilde u\,
\Gamma^*\right)\epsilon_\ell
                 +\left(v+\tilde v\,\Gamma^*\right)\epsilon_\rr,
\end{equation}

                其中
% 式 (5.15)
\setcounter{equation}{14}
\begin{equation}\label{eq:ujbyty}
\Gamma^*=\begin{cases}\Gamma_{01}&\text{in Lorentz
signature}\\ i\Gamma_{01}&\text{in Euclidean signature}\\\end{cases}
\end{equation}
 是区分二维旋量两种手征的算子。这里我们将采用欧几里得号差，
因为这对拓扑场论更自然。

\eqref{eq:blop}\ 中的超对称生成元恰好就是满足下列条件的那些：
% 式 (5.16)
\setcounter{equation}{15}
\begin{equation}\label{eq:zonop}
\begin{aligned}\Gamma_{2367}\epsilon& = \epsilon \\
\Gamma_{0145}\epsilon & = \epsilon.\\\end{aligned}
\end{equation}
 第一个是式 \eqref{eq:ytytag}，
它说由 $\epsilon$ 生成的超对称不被 $C$ 的曲率破缺；第二个说它
同样不被 $\Sigma$ 的曲率破缺。有了这个更一般的出发点，就有可能从
四维规范理论得到二维 TQFT 的整个族
$\Bbb{CP}_h^1\times \Bbb{CP}_h^1$，尽管不是从四维
TQFT 得到。

实际上，把 \eqref{eq:blop}\ 改写在 $\Gamma^*$ 的本征基中是很方便的：
% 式 (5.17)
\setcounter{equation}{16}
\begin{equation}\label{eq:bilop}
\hat\epsilon=\half\left(1-\Gamma^*\right)(u'\epsilon_\ell+v'\epsilon_r)
+\half\left(1+\Gamma^*\right)(u''\epsilon_\ell+v''\epsilon_r).
\end{equation}

这里 $(u',v')$ 与 $(u'',v'')$ 分别是 $\Bbb{CP}^1_h\times
\Bbb{CP}^1_h$ 两个因子的齐次坐标。设 $\hat\delta_T$ 是由 $\hat \epsilon$ 生成的扩充拓扑
对称。为确定两个复结构 $(J_+,J_-)$，或等价地确定对
$(w_+,w_-)$，我们只需计算
$\hat\delta_T\chi^+_{\bar y\,\bar z}$ 与
$\hat\delta_T\chi^-_{y\bar z}$。把它们置零就给出
\eqref{eq:onso} 的推广。为确定式 \eqref{eq:helpme} 的通常公式的推广，我们只需要知道：
若 $\tilde\chi^+=\Gamma^{\bar y\,\bar z}\chi^+_{\bar y\,\bar z}$
且 $\tilde\chi^-=\Gamma^{y\bar z}\chi^-_{y\bar z}$，则
$\tilde\chi^+\Gamma^*=\tilde\chi^+$ 且
$\tilde\chi^-\Gamma^*=-\tilde\chi^-$。利用这一点，我们得到
% 式 (5.18)
\setcounter{equation}{17}
\begin{equation}\label{eq:qengo}
\begin{aligned}\hat\delta_T\chi^+_{\bar y\,\bar
z}&=u'(F-\phi\wedge\phi)_{\bar y\,\bar z}+v'(D\phi)_{\bar y\,\bar
z}\\ \hat\delta_T\chi^-_{ y \bar z}&=v''(F-\phi\wedge\phi)_{y
z}-u''(D\phi)_{yz}.\\\end{aligned}
\end{equation}
 于是，若 $t'=v'/u'$ 且
$t''=v''/u''$，则式 \eqref{eq:igot}\ 与
\eqref{eq:nigot}\ 的推广为 % 式 (5.19)
\setcounter{equation}{18}
\begin{equation}\label{eq:pokodo}
\begin{aligned}\partial_{\bar y}(A_{\bar
z}+t'\phi_{\bar z})&=0\\ \partial_y(A_{\bar
z}-(t'')^{-1}\phi_{\bar z})&=0 .\\\end{aligned}
\end{equation}
 这确定了两个
复结构： % 式 (5.20)
\setcounter{equation}{19}
\begin{equation}\label{eq:plombob}
(w_+,w_-)=(-t',(t'')^{-1}).
\end{equation}

现在我们可以确定 $S$-对偶如何作用在整个族
$\Bbb{CP}^1_h\times \Bbb{CP}^1_h$ 上。我们在式 \eqref{eq:duco}\ 中了解到，
对偶变换 $\left(\begin{matrix}a & b \\ c & d\end{matrix}\right)\in
SL(2,\Bbb{Z})$（或它对非单连通型规范群的类比）
作用在 $t$ 上为 % 式 (5.21)
\setcounter{equation}{20}
\begin{equation}\label{eq:utlo}
t\to t{c\tau+d\over
|c\tau+d|}.
\end{equation}
  同样的推理说明 $t'$ 与 $t''$
以完全相同的方式变换。由此及 \eqref{eq:plombob} 可知，$S$-对偶对 $\Bbb{CP}^1_h\times \Bbb{CP}^1_h$ 的作用是
\begin{equation*}
\tag*{(5.22)}
\begin{aligned}w_+ & \to w_+{c \tau+d\over
|c\tau+d|}\\
                   w_-  & \to w_-{|c\tau+d|\over c
                   \tau+d}.\\\end{aligned}
\end{equation*}

一个重要的特例是，复结构 $I$ 中的 $B$-模型（对应于 $(w_+,w_-)=(0,0)$）在
对偶变换下完全不变。复结构 $I$ 中的 $A$-模型（对应于
$(w_+,w_-)=(0,\infty)$）同样不变。关于二维 TQFT 的这一结论以及其他不在
特殊族 $\Bbb{CP}^1_g$ 上的结论将是有用的辅助
工具，但正如将要变得清楚的，几何 Langlands 纲领
实际上是关于那个特殊族的一个陈述或一族陈述。

\par\medskip\emph{对度量的依赖}

一般说来，靶空间为 $\MH(G,C)$ 的二维 TQFT 将
依赖于 $C$ 的复结构，因为这会影响
$\MH(G,C)$ 的超凯勒结构。然而，若一个二维 TQFT（例如由
$\Bbb{CP}^1_g$ 参数化的那些）来自四维理论，它就必须
与 $C$ 上的复结构无关。毕竟，$M=\Sigma\times C$ 上的
四维 TQFT 与 $M$ 上的度量无关，特别地与 $C$ 上的度量无关。

在某些情形下，我们可以直接看出这些模型与 $C$ 的复结构无关。在 $t=\pm i$ 时，我们
得到复结构 $\pm J$ 中的 $B$-模型；正如我们在
第~\ref{sec:mhhyper}~节中观察到的，这个复结构与 $C$ 的度量无关。对实的 $t$，我们得到某个复结构中的 $A$-模型，
该复结构是复结构 $I$ 与 $K$ 的线性组合。$A$-模型由相应的
辛结构决定，而该辛结构是辛结构 $\omega_I$ 与 $\omega_K$ 的线性组合；我们在第~\ref{sec:mhhyper}~节中发现它们与 $C$ 的度量无关。

为了走得更远，我们将使用广义复几何。

\begin{yizhu}[原文疑误]
本片「第一个是式 \eqref{eq:ytytag}」处，原文指 §5 的式 (5.3) $\Gamma_{2367}\epsilon=\epsilon$（不是 §4 的 (4.24)）；harvmac 源码对该公式只给了 \texttt{tag*}、未给编号标签，译本已补标签 \texttt{eq:ytytag}，引用正确。按：这是排版脚手架问题，不是原文笔误。
\end{yizhu}

\begin{yizhu}[原文疑误]
原文「a few simple simple observations」中 simple 重复了一次，显系笔误（应为一个 simple）；译文按「几点简单的观察」译出。
\end{yizhu}

\begin{yizhu}[原文疑误]
原文「this does not meant that they cannot be described」中 meant 应为 mean，属笔误；译文按「并不意味着」译出。
\end{yizhu}

\begin{yizhu}[易错点]
$w_+=-t$、$w_-=t^{-1}$ 这一嵌入是在把 $G$ 阿贝尔化（破缺到极大环面）之后算出的；它给出四维扭化参数 $t$（即 $\Bbb{CP}^1_g$ 上的仿射参数）与二维复结构对 $(J_+,J_-)$ 的对应关系。注意 $t$ 与 $-t$ 给出二维中相互等价的理论，这来自 $SU(4)_{\cal R}$ 的中心，而不是来自四维层面的等价。
\end{yizhu}

\begin{yizhu}[含糊点拨]
$B$-模型要求 $w_+=w_-$，即 $t^{-1}=-t$，故只有 $t=\pm i$；$A$-模型要求两个复结构相反，而 $w\to-1/\bar w$ 把 $I_w$ 映到相反复结构，故 $t$ 为实数即可（$t=0,\infty$ 也包含在内）。按：这两条是本片后面「$B$-模型不变、$A$-模型在实 $t$ 处出现」的几何根据，也是把四维 $S$-对偶翻译成二维镜像对称的入口。
\end{yizhu}

\begin{yizhu}[物理图像]
点 $(0,0)$ 对应于复结构 $I$ 中的 $B$-模型，它是数学上几何 Langlands 的出发点；Hitchin 纤维化纤维上的 $T$-对偶（在其上作用为 Fourier-Mukai 变换）把该 $B$-模型映到自身。虽然 $(0,0)$ 不在来自四维 TQFT 的族 $\Bbb{CP}^1_g$ 中，但把它取无穷小扰动 $(\alpha,\beta)$ 就得到 $\Bbb{CP}^1_g$ 上任意点附近的 $\Bbb{C}^*$ 轨道；原文把它作为一个可能的将来方向提出。
\end{yizhu}

\begin{yizhu}[易错点]
「复结构 $\hat J$ 中的 $A$-模型」只是简称：原文明确指出 $A$-模型真正依赖的是辛形式 $\omega_{\hat J}$，而不是复结构本身。后文把四维 $S$-对偶化归为二维镜像对称时，$A$-模型这一侧对应的正是 $\MH$ 的辛结构 $\omega_K$。
\end{yizhu}
